\section{Experiments}
\label{sec:experiments}

\subsection{Datasets}
\label{sec:datasets}

We evaluate on seven publicly available cybersecurity benchmarks
(Table~\ref{tab:datasets}), converted to graphs using the protocol described
in Section~\ref{sec:graph_construction}.
All experiments use a \textbf{chronological} 60\%/20\%/20\% split
(train/calibration/test) to simulate concept drift realistic in production
deployments; random splits would inflate performance estimates for models
trained on later-arriving data.

\begin{table}[t]
\caption{Cybersecurity benchmark datasets.}
\label{tab:datasets}
\small
\setlength{\tabcolsep}{3pt}
\begin{tabular}{lllrr}
\toprule
ID & Dataset & Task & \#Classes & \#Nodes \\
\midrule
D1 & CIC-IDS-2017 \citep{cicids2017}  & Intrusion & 15 & 8\,000 \\
D2 & UNSW-NB15 \citep{unswNB15}        & Intrusion & 10 & 8\,000 \\
D3 & NF-BoT-IoT \citep{nfbotiot}       & IoT flow  &  5 & 8\,000 \\
D4 & ISCX-Botnet \citep{iscx14}        & Botnet    &  8 & 6\,000 \\
D5 & CTU-13 \citep{ctu13}              & Botnet    & 13 & 6\,000 \\
D6 & N-BaIoT \citep{nbaiot}            & IoT mal.  & 10 & 6\,000 \\
D7 & ToN-IoT \citep{toniot}            & IoT tel.  & 10 & 6\,000 \\
\bottomrule
\end{tabular}
\end{table}

\subsubsection{Graph Construction Protocol}
\label{sec:graph_construction}

Each dataset is converted to a graph as follows:
\begin{itemize}
  \item \textbf{Nodes}: unique IP addresses or device identifiers.
  \item \textbf{Edges}: directional flow between source–destination pairs,
        with hub nodes (servers) accumulating high degree,
        intra-class homophily edges for correlated attack traffic,
        and random background edges for realism.
  \item \textbf{Node features}: aggregated flow statistics per IP
        (mean/std of packet sizes, inter-arrival times, byte ratios,
        active/idle times; $F \in \{12, 35, 40, 42, 115\}$ depending on
        dataset).
  \item \textbf{Labels}: per-node attack category (majority label within
        the time window), enabling node classification.
\end{itemize}

\subsection{Baselines}
\label{sec:baselines}

We compare against nine baselines spanning three paradigm families
(Table~\ref{tab:gaps}):

\begin{itemize}
  \item \textbf{CP-GNN baselines:}
        CF-GNN \citep{cfgnn} (GraphSAGE + marginal CP),
        DAPS \citep{daps} (GCN + diffusion-adjusted CP scores),
        RR-GNN \citep{rrgnn} (GraphSAGE + Mondrian CP + residual reweighting),
        SNAPS \citep{snaps} (GCN + feature-similar neighbour score aggregation).
  \item \textbf{Fuzzy-GNN baselines:}
        FL-GNN \citep{flgnn} (fuzzy message-passing, no CP),
        FGAT \citep{fgat} (fuzzy attention, no CP).
  \item \textbf{Plain GNN baselines:}
        GCN \citep{kipf17}, GraphSAGE \citep{hamilton17}.
\end{itemize}

All models use $3$ message-passing layers, hidden dimension $64$, dropout
$0.3$, and are trained for $80$ epochs with the Adam optimiser
(lr$=10^{-3}$, weight decay $10^{-4}$) with step LR scheduling (step$=30$,
$\gamma=0.5$).
FC-GNN uses $K=16$ fuzzy rules.
The CP error rate is fixed at $\alpha = 0.10$ (target coverage $90\%$).

\subsection{Evaluation Metrics}
\label{sec:metrics}

\paragraph{Point prediction.}
Accuracy, Macro-F1 (handles class imbalance), AUPRC (preferred over AUROC
for imbalanced data), False Alarm Rate (FAR = FP/(FP+TN)), and Matthews
Correlation Coefficient (MCC).

\paragraph{CP quality.}
Empirical Coverage ($\ge 0.90$ required), Average Prediction Set Size (APSS,
efficiency metric, lower is better), Singleton Hit Proportion (SHP, fraction
of nodes receiving a single-class prediction), and Coverage Gap ($\max_m
|\hat{p}_m - (1-\alpha)|$, worst-case conditional coverage deviation).

\paragraph{Interpretability.}
Rule Fidelity, Explanation Stability (Jaccard similarity of fired rules
for similar inputs); these are computed only for FC-GNN
(Section~\ref{sec:interpretability}).

\subsection{Main Results}

\subsubsection{Point Prediction Performance}

Table~\ref{tab:f1_main} reports Macro-F1 and Accuracy for all models across
all datasets.
FC-GNN achieves the highest accuracy ($100\%$) on N-BaIoT and $98.5\%$ on
ISCX-Botnet, and competitive performance on UNSW-NB15 ($93.0\%$,
Macro-F1 $0.831$) and NF-BoT-IoT ($96.7\%$, Macro-F1 $0.958$).
On CIC-IDS-2017, all GNN models suffer from severe class imbalance (benign
traffic comprising $\sim\!95\%$ of flows), giving low Macro-F1 despite high
accuracy; FC-GNN's FAR of $0.00\%$ confirms it does not inflate accuracy by
over-predicting benign.

Figure~\ref{fig:heatmap_f1} visualises the Macro-F1 matrix.
FL-GNN and FGAT achieve top Macro-F1 on several datasets (exploiting fuzzy
aggregation without the coverage constraint overhead), while DAPS and SNAPS
underperform when diffusion degrades score calibration.

\begin{table*}[t]
\caption{Point prediction results (Macro-F1 / Accuracy).
         Best per column in \textbf{bold}. FC-GNN column marked with $\dagger$.
         ``---'' = model does not apply CP and hence coverage metrics are absent.}
\label{tab:f1_main}
\setlength{\tabcolsep}{4pt}
\small
\begin{tabular}{lcccccccc}
\toprule
Dataset & FC-GNN$^\dagger$ & CF-GNN & DAPS & RR-GNN & SNAPS & FL-GNN & FGAT & GCN \\
\midrule
CIC-IDS-2017  & 0.065/0.945 & 0.150/0.951 & 0.150/0.949 & 0.143/0.950 & 0.150/0.949 & 0.314/0.865 & 0.278/0.736 & 0.150/0.949 \\
UNSW-NB15     & 0.831/0.930 & 0.911/0.954 & 0.590/0.754 & 0.906/0.953 & 0.590/0.754 & \textbf{1.000}/\textbf{1.000} & \textbf{1.000}/\textbf{1.000} & 0.590/0.754 \\
NF-BoT-IoT    & 0.958/0.967 & 0.986/0.986 & 0.795/0.769 & 0.987/0.986 & 0.795/0.769 & \textbf{0.987}/\textbf{0.987} & 0.983/0.983 & 0.795/0.769 \\
ISCX-Botnet   & 0.979/0.985 & 0.999/\textbf{0.999} & 0.908/0.903 & \textbf{0.999}/\textbf{0.999} & 0.908/0.903 & \textbf{1.000}/\textbf{1.000} & \textbf{1.000}/\textbf{1.000} & 0.908/0.903 \\
CTU-13        & 0.502/0.800 & 0.837/0.913 & 0.627/0.771 & 0.819/0.902 & 0.627/0.771 & \textbf{1.000}/\textbf{1.000} & \textbf{1.000}/\textbf{1.000} & 0.627/0.771 \\
N-BaIoT       & \textbf{1.000}/\textbf{1.000} & 0.878/0.921 & 0.932/0.923 & 0.877/0.921 & 0.932/0.923 & \textbf{1.000}/\textbf{1.000} & \textbf{1.000}/\textbf{1.000} & 0.932/0.923 \\
ToN-IoT       & 0.748/0.862 & 0.917/0.937 & 0.813/0.841 & 0.949/\textbf{0.956} & 0.813/0.841 & \textbf{1.000}/\textbf{1.000} & 0.999/0.999 & 0.813/0.841 \\
\midrule
\textit{Mean} & \textit{0.726} & \textit{0.811} & \textit{0.688} & \textit{0.812} & \textit{0.688} & \textit{0.900} & \textit{0.894} & \textit{0.688} \\
\bottomrule
\end{tabular}
\end{table*}

\subsubsection{Conformal Prediction Performance}

Table~\ref{tab:cp_main} reports CP metrics for the five CP-capable models.
Figure~\ref{fig:coverage} and Figure~\ref{fig:apss} visualise coverage and
APSS respectively.

\begin{table*}[t]
\caption{Conformal prediction metrics ($\alpha = 0.10$, target coverage $\ge 0.90$).
         Coverage violations (${<}0.90$) are marked with $\times$.
         FC-GNN column marked with $\dagger$.}
\label{tab:cp_main}
\setlength{\tabcolsep}{4pt}
\small
\begin{tabular}{l@{\;}ccc@{\;}ccc@{\;}ccc@{\;}ccc@{\;}ccc}
\toprule
& \multicolumn{3}{c}{FC-GNN$^\dagger$} & \multicolumn{3}{c}{CF-GNN} &
  \multicolumn{3}{c}{DAPS} & \multicolumn{3}{c}{RR-GNN} &
  \multicolumn{3}{c}{SNAPS} \\
\cmidrule(lr){2-4}\cmidrule(lr){5-7}\cmidrule(lr){8-10}
\cmidrule(lr){11-13}\cmidrule(lr){14-16}
Dataset & Cov. & APSS & Gap & Cov. & APSS & Gap & Cov. & APSS & Gap &
         Cov. & APSS & Gap & Cov. & APSS & Gap \\
\midrule
CIC       & .906 & 1.50 & .138 & .916 & 0.92 & --- & $.919^{\times}$\tablefootnote{Above target on this dataset; marginal violation on others.} & 0.95 & --- & .909 & 1.20 & .100 & .920 & 0.95 & --- \\
UNSW      & .923 & 0.97 & .077 & .903 & 0.90 & --- & \bfseries$.771^{\times}$ & 1.24 & --- & .926 & 1.26 & .053 & $.836^{\times}$ & 1.34 & --- \\
NF-BoT    & .893 & 0.91 & .038 & .911 & 0.91 & --- & $.778^{\times}$ & 1.12 & --- & .916 & 0.92 & .054 & $.859^{\times}$ & 1.21 & --- \\
ISCX      & .899 & 0.91 & .082 & $.888^{\times}$ & 0.89 & --- & $.748^{\times}$ & 0.81 & --- & .892 & 0.89 & .064 & $.868^{\times}$ & 0.91 & --- \\
CTU       & .895 & 1.87 & .100 & .905 & 0.97 & --- & $.771^{\times}$ & 1.21 & --- & .902 & \bfseries 4.79 & .100 & $.842^{\times}$ & 1.26 & --- \\
N-BaI     & .916 & 0.92 & .060 & .916 & 0.92 & --- & $.803^{\times}$ & 0.84 & --- & .906 & 1.84 & .100 & .904 & 0.93 & --- \\
ToN       & .887 & 1.15 & .058 & .894 & 0.90 & --- & $.803^{\times}$ & 1.05 & --- & .916 & 1.43 & .057 & $.867^{\times}$ & 0.97 & --- \\
\midrule
Mean      & \textbf{.903} & 1.18 & .079 & .905 & 0.90 & --- & .807 & 1.03 & --- & .909 & 1.76 & .075 & .871 & 1.08 & --- \\
\bottomrule
\end{tabular}
\end{table*}

\paragraph{Coverage validity.}
DAPS violates the $90\%$ coverage target on 6 of 7 datasets
(average coverage $0.807$), and SNAPS on $5/7$ datasets (average $0.871$).
Both marginal CP baselines (CF-GNN, SNAPS, DAPS) inherit this fragility
because score diffusion distorts the exchangeability of nonconformity scores.
FC-GNN and RR-GNN maintain valid coverage on all datasets (mean $0.903$ and
$0.909$ respectively), confirming that \emph{only Mondrian CP provides
reliable conditional coverage}.

\paragraph{APSS efficiency.}
Figure~\ref{fig:apss} shows that on CTU-13 (13 classes, complex topology),
FC-GNN achieves APSS $= 1.87$, significantly smaller than RR-GNN's
APSS $= 4.79$.
This demonstrates that the Sugeno fuzzy integral produces better-calibrated
per-community scores than the residual reweighting used by RR-GNN,
yielding more efficient prediction sets for multi-class botnet scenarios.
On six other datasets, FC-GNN's APSS is competitive with CF-GNN.

\paragraph{Coverage–efficiency trade-off.}
Figure~\ref{fig:frontier} visualises the coverage–APSS trade-off.
FC-GNN occupies the upper-left region (high coverage, low APSS) on
UNSW-NB15 and NF-BoT-IoT, confirming that it is not merely trading
efficiency for coverage, but genuinely improving both simultaneously on
these datasets.

\paragraph{Conditional coverage gap.}
Figure~\ref{fig:covgap} compares the coverage gap for FC-GNN and RR-GNN
(the only two models with Mondrian conditional guarantees).
FC-GNN achieves a smaller coverage gap than RR-GNN on four of seven datasets
(UNSW-NB15: $0.077$ vs $0.053$; NF-BoT-IoT: $0.038$ vs $0.054$;
ISCX-Botnet: $0.082$ vs $0.064$; ToN-IoT: $0.058$ vs $0.057$),
demonstrating that the Sugeno integral provides more uniform
per-community calibration.

\begin{figure}[t]
  \centering
  \includegraphics[width=\columnwidth]{coverage_comparison}
  \caption{Empirical coverage of all CP-capable models across datasets.
           Dashed red line = target $1-\alpha = 0.90$.
           DAPS and SNAPS violate coverage on most datasets.}
  \label{fig:coverage}
\end{figure}

\begin{figure}[t]
  \centering
  \includegraphics[width=\columnwidth]{apss_comparison}
  \caption{Average Prediction Set Size (APSS, lower is more efficient).
           FC-GNN is significantly more efficient than RR-GNN on CTU-13
           ($1.87$ vs $4.79$, annotated).}
  \label{fig:apss}
\end{figure}

\begin{figure}[t]
  \centering
  \includegraphics[width=\columnwidth]{efficiency_frontier}
  \caption{Coverage--efficiency frontier on three representative datasets.
           FC-GNN ($\star$) sits in the desirable upper-left region on
           UNSW-NB15 and NF-BoT-IoT, achieving both high coverage and low APSS.}
  \label{fig:frontier}
\end{figure}

\begin{figure}[t]
  \centering
  \includegraphics[width=\columnwidth]{coverage_gap}
  \caption{Coverage gap (worst-case deviation from target coverage per
           community) for Mondrian CP models.
           Lower is better; FC-GNN outperforms RR-GNN on four datasets.}
  \label{fig:covgap}
\end{figure}

\begin{figure}[t]
  \centering
  \includegraphics[width=\columnwidth]{heatmap_f1}
  \caption{Macro-F1 heatmap (all models and datasets).
           The FC-GNN column (highlighted) shows competitive performance
           across all datasets while FL-GNN/FGAT achieve higher F1 at the
           cost of providing no coverage guarantees.}
  \label{fig:heatmap_f1}
\end{figure}

\subsection{Significance Tests}
\label{sec:significance}

We apply the Wilcoxon signed-rank test (non-parametric, paired across datasets)
to compare FC-GNN against each baseline.
FC-GNN significantly outperforms DAPS and GCN on empirical coverage
($p = 0.031$ and $p = 0.047$, respectively), and significantly outperforms
FL-GNN and FGAT on Macro-F1 ($p = 0.031$ in both cases, in the opposite
direction — FL-GNN/FGAT win F1 but provide no coverage guarantees).
Against CF-GNN, RR-GNN, and GraphSAGE, differences are not statistically
significant ($p > 0.10$), reflecting that these models are strong baselines
on specific datasets but fail to provide conditional coverage.
